






\hypersetup{colorlinks=true,linkcolor=blue,citecolor=blue}




















\title[Short title]{On 1111111111111 rrry regularity of weak solutions to a class of elliptic equations}
\author{Yang Jianwei} 


\thanks{This work was partially supported by NSFC (Grant No. 00000000).}

\begin{document}

\begin{abstract}ffffgfgffgffgf ghdufhujvijijijvij hhjfffujhhhh h  jjjj
\end{abstract}

\subjclass[2020]{Primary 35J60; Secondary 35B65, 46E35}
\keywords{partial differential equations, Sobolev space, weak solution, regularity}
\maketitle

\section{Introduction}\label{sec:intro}
For the classical background and related results, see~\cite{GT,LU}.

Let $\Om\subset\R^n$ ($n\ge 2$) be a bounded domain with Lipschitz boundary.
We consider the quasilinear elliptic problem
\begin{equation}\label{eq:main}
  -\operatorname{div}\bigl(A(x,\nabla u)\bigr)=f \quad \text{in } \Om,
  \qquad u=0 \quad \text{on } \partial\Om ,
\end{equation}
where $A:\Om\times\R^n\to\R^n$ satisfies the standard $p$-growth and
coercivity conditions with $1<p<\infty$.

\begin{definition}\label{def:weak}
A function $u\in W^{1,p}_0(\Om)$ is called a \emph{weak solution} of
\eqref{eq:main} if for every $\varphi\in W^{1,p}_0(\Om)$,
\begin{equation}\label{eq:weak}
  \int_{\Om} A(x,\nabla u)\cdot\nabla\varphi\,dx
  = \int_{\Om} f\varphi\,dx .
\end{equation}
\end{definition}

\begin{theorem}[Main result]\label{thm:main}
Assume $f\in L^q(\Om)$ with $q>n/p$ and let $u\in W^{1,p}_0(\Om)$ be a weak
solution of \eqref{eq:main}. Then there exist $\alpha\in(0,1)$ and $C>0$,
depending only on $n,p,q$, and the structural constants, such that
\begin{equation}\label{eq:holder}
  \|\nabla u\|_{C^{0,\alpha}(\Om')} \le C
  \quad\text{for every } \Om'\Subset\Om .
\end{equation}
\end{theorem}

\subsection{Strategy of the proof}
The proof combines a Caccioppoli inequality with an iteration argument.
We first derive the energy estimate
\begin{equation}
  \int_{B_r} |\nabla u|^p\,dx
  \le C r^{-p}\int_{B_{2r}} |u-u_{2r}|^p\,dx + C\int_{B_{2r}}|f|^{p'}\,dx ,
\end{equation}
and then apply the De Giorgi lemma on nested balls; see
Section~\ref{sec:proof}.

\section{Preliminaries}\label{sec:pre}

We recall the Sobolev embedding $W^{1,p}_0(\Om)\hookrightarrow L^{p^*}(\Om)$
with $p^*=np/(n-p)$ for $p<n$, and the Poincaré inequality
$\|u\|_{L^p}\le C(\Om)\|\nabla u\|_{L^p}$.

\begin{lemma}[Caccioppoli]\label{lem:cacc}
For any $B_{2r}\Subset\Om$ and any $k\in\R$,
\begin{equation}
  \int_{B_r\cap\{u>k\}}|\nabla u|^p\,dx
  \le \frac{C}{r^p}\int_{B_{2r}\cap\{u>k\}}(u-k)^p\,dx .
\end{equation}
\end{lemma}

\section{Proof of Theorem~\ref{thm:main}}\label{sec:proof}

\begin{proof}[Proof]
Testing \eqref{eq:weak} with $\varphi=\eta^p(u-k)_+$ and applying
Lemma~\ref{lem:cacc} yields a decay estimate for the level sets.
Iteration gives an $L^\infty$ bound on $\nabla u$, and a standard
rescaling argument then implies \eqref{eq:holder}.
\end{proof}

\begin{remark}
The assumption $q>n/p$ is sharp in the scale of Lebesgue spaces;
counterexamples are known for $q=n/p$.
\end{remark}

\section*{Acknowledgments}
The author thanks the referee for careful reading and valuable suggestions.

\begin{thebibliography}{9}
\bibitem{GT} D.~Gilbarg and N.~S.~Trudinger, \emph{Elliptic Partial Differential
  Equations of Second Order}, 2nd ed., Springer, 2001.
\bibitem{LU} O.~A.~Ladyzhenskaya and N.~N.~Ural'tseva, \emph{Linear and Quasilinear
  Elliptic Equations}, Academic Press, 1968.
\end{thebibliography}

\end{document}
